\documentclass[11pt,a4paper]{article}

% ----------------------------------------------------------------------
% Protocol: raw recording -> BuzzSuite app (intake, tracking, plot)
% Written for someone with very little coding experience.
% ----------------------------------------------------------------------

\usepackage[utf8]{inputenc}
\usepackage[T1]{fontenc}
\usepackage[margin=2.5cm]{geometry}
\usepackage{enumitem}
\usepackage{xcolor}
\usepackage{tcolorbox}
\usepackage{graphicx}
\usepackage{hyperref}
\usepackage{titlesec}

\definecolor{appblue}{RGB}{30,90,160}
\definecolor{lightgray}{RGB}{240,240,240}
\definecolor{treegray}{RGB}{250,250,245}

\hypersetup{
  colorlinks=true,
  linkcolor=appblue,
  urlcolor=appblue
}

\titleformat{\section}{\Large\bfseries\color{appblue}}{\thesection.}{0.5em}{}
\titleformat{\subsection}{\large\bfseries}{\thesubsection}{0.5em}{}

% A box for important tips / warnings
\newtcolorbox{tipbox}[1]{colback=lightgray,colframe=appblue,
  fonttitle=\bfseries,title=#1,boxrule=0.6pt,arc=2mm}

% A box for the folder-tree diagram
\newtcolorbox{treebox}{colback=treegray,colframe=appblue!60,
  boxrule=0.6pt,arc=2mm,left=8pt,right=8pt,top=6pt,bottom=6pt}

\setlength{\parindent}{0pt}
\setlength{\parskip}{6pt}

\title{\textbf{BuzzSuite: From Recording to Data}\\[4pt]
\large A Simple Step-by-Step Protocol}
\author{}
\date{}

\begin{document}
\maketitle

\vspace{-1.5em}
\begin{center}
\textit{No coding required --- just clicking, naming things, and selecting files.}\\[2pt]
\textit{\small The Activity module used in this protocol is BuzzSuite's activity-tracking
pipeline (built on the original BuzzWatch tracking engine) --- this protocol covers that
path from raw video to the activity plot, start to finish.}
\end{center}

\hrule
\vspace{1em}

\begin{tipbox}{The workflow at a glance}
\begin{enumerate}[leftmargin=*,itemsep=1pt]
  \item Collect your \texttt{.h264} video data from the camera.
  \item Open \textbf{BuzzSuite} and use the \textbf{Intake wizard} to load the video and create
        the experiment --- this also converts \texttt{.h264} to \texttt{.mp4} for you, no separate
        app needed.
  \item Extract a background image and draw the cage borders.
  \item Run tracking.
  \item Make the activity plot.
\end{enumerate}
\end{tipbox}

% ----------------------------------------------------------------------
\section{Step 1 --- Collect your \texttt{.h264} data}

Copy the \texttt{.h264} video files off the \textbf{Raspberry Pi} onto the
computer. Put them somewhere you can find them (you will point BuzzSuite at
them in Step~2).

% ----------------------------------------------------------------------
\section{Step 2 --- Open BuzzSuite and create the experiment}

\subsection*{2.1 --- Launch the app}

Double-click \texttt{run\_buzzsuite.bat} (or its desktop shortcut) in the BuzzSuite folder.

\begin{tipbox}{If Windows shows a blue ``Windows protected your PC'' warning}
Click \textbf{``More info''}, then \textbf{``Run anyway''}. This is normal and
safe here.
\end{tipbox}

A welcome screen appears first. Click \textbf{``New experiment $\rightarrow$''} to jump straight
into setup (or, if you're continuing work on an experiment you already created, select it from the
\textbf{recent experiments} list, or use \textbf{``Browse for experiment\ldots''}, then click
\textbf{``Continue $\rightarrow$''}).

\subsection*{2.2 --- The Intake wizard}
\textit{(Tab: \texttt{1} \textperiodcentered\ \texttt{Setup} $\to$ sub-tab \texttt{Experiment \& Video})}

This one screen replaces what used to be a separate folder-creation step and a separate Convert
app. Work through it top to bottom:

\begin{enumerate}[label=\textbf{\Alph*\ ---}\ ,leftmargin=*,itemsep=8pt]
  \item \textbf{Source video folder.} Click \textbf{Browse\ldots} and select the folder
        with your \texttt{.h264} files (from Step~1), then click \textbf{Scan}. BuzzSuite lists
        how many video segments and dates it found.

  \item \textbf{Dates to include.} A checkbox appears for every recording date found. All
        are ticked by default --- untick any date you don't want included.

  \item \textbf{Experiment setup.} Fill in these boxes, then click
        \textbf{``Create / open experiment''}:

\begin{description}[leftmargin=0pt,style=nextline]
  \item[Experiment:] type \texttt{Flare}. (Empty by default.)
  \item[Incubator:] \texttt{Cakung} \quad or \quad \texttt{Bangkok} \quad (pick from the dropdown)
  \item[Species:] e.g.\ \texttt{AedesAegypti}
  \item[Sex:] \texttt{F} or \texttt{M} (pick from the dropdown)
  \item[Cage name:] BuzzSuite fills this in for you from Species + Sex + the dates you ticked in
    Step~B (e.g.\ \texttt{AedesAegypti\_M\_20260721\_20260722}). \textbf{You can type over it}
    if you'd rather keep the older short style, e.g.\ \texttt{AedesAegyptiM\_21-22}.
\end{description}

\begin{tipbox}{Don't mix up the two ``Cage'' fields}
\textbf{Incubator} (\texttt{Cakung}/\texttt{Bangkok}) is what used to be called ``Cage Name''.
\textbf{Cage name} is a different box --- it's the actual experiment folder name (what used to be
called ``Batch Name''). Same idea as before, just split across two clearly-named boxes now.
\end{tipbox}

BuzzSuite creates the experiment folder for you automatically --- there is no more manual folder
creation in File Explorer. For reference, this is what gets built behind the scenes:

\begin{treebox}
\ttfamily
\textless your data drive\textgreater\textbackslash Buzzwatch\\
\hspace*{1em}\textbackslash--\ Recording\\
\hspace*{3em}\textbackslash--\ Flare\hfill\textnormal{\itshape\small (= your ``Experiment'' box)}\\
\hspace*{5em}\textbackslash--\ Cakung\hfill\textnormal{\itshape\small (= your ``Incubator'' box)}\\
\hspace*{7em}\textbackslash--\ \textbf{AedesAegyptiM\_21-22}\hfill\textnormal{\itshape\small (= your ``Cage name'' box)}\\
\end{treebox}

You don't need to know or type this path --- BuzzSuite finds your data drive automatically.

  \item \textbf{Convert / move video into the experiment.} Click
        \textbf{``Convert / move selected dates''}. Any \texttt{.h264} files are losslessly
        converted to \texttt{.mp4}; any files that are already \texttt{.mp4} are copied straight
        in. A progress bar shows how far along it is. Wait for it to finish.
\end{enumerate}

\begin{tipbox}{A separate H264 $\to$ MP4 Converter tab also exists}
If you ever want to convert video without going through the wizard (e.g.\ into a folder outside a
BuzzSuite experiment), there's also a standalone \textbf{``H264 $\to$ MP4 Converter''} tab along
the top of the app that does the same conversion. For the normal workflow, Step~D above is all you
need.
\end{tipbox}

% ----------------------------------------------------------------------
\section{Step 3 --- Background image and cage borders}
\textit{(Tab: \texttt{1} \textperiodcentered\ \texttt{Setup} $\to$ sub-tab \texttt{Background \& Cage Border})}

\begin{enumerate}[leftmargin=*]
  \item \textbf{Extract Images from Video.}\\
        Click the button and wait until it finishes.

  \item \textbf{Get Background from Images.}\\
        Click the button and wait. \emph{This one takes a while to run.}

  \item \textbf{Draw the cage borders.}
  \begin{enumerate}[label=\alph*.]
    \item Click \textbf{Show background with borders} (the cage image appears on the right).
    \item Click \textbf{Draw Cage Borders} --- a separate image window opens.
    \item Click the \textbf{four corners of the cage} --- \textbf{any order is fine},
          BuzzSuite automatically works out which corner is which.
    \item After the 4th click, a box is drawn joining the corners. Press any key to close the
          window --- the borders are saved automatically.
  \end{enumerate}
\end{enumerate}

\begin{center}
\begin{verbatim}
        (*)---------------(*)
         |                 |
         |      cage       |
         |                 |
        (*)---------------(*)
\end{verbatim}
\textit{\small click the four corners in any order}
\end{center}

If you ever need to redo this, click \textbf{Reset Drawn Borders} first.

% ----------------------------------------------------------------------
\section{Step 4 --- Run tracking}

There are two ways to do this --- pick whichever you prefer. \textbf{4.2 (the Experiment
Dashboard) works perfectly well for a single experiment too}, not just several at once, so
don't feel you need to use 4.1 just because you're only tracking one.

\subsection*{4.1 --- The Single experiment tab}
\textit{(Tab: \texttt{2} \textperiodcentered\ \texttt{Analysis} $\to$ sub-tab \texttt{Single experiment})}

\begin{enumerate}[leftmargin=*]
  \item Choose \textbf{what to track}:
  \begin{itemize}
    \item \textbf{All untracked (batch)} --- the normal choice; tracks every video segment that
          hasn't been tracked yet and skips ones already done. \emph{(This is selected by
          default.)}
    \item \textbf{Specific segments} --- pick individual segments from a list.
    \item \textbf{Time range} --- pick a start and end time.
  \end{itemize}

  \item Click \textbf{``Run tracking''}. \emph{This normally takes several hours} --- you don't
        need to set anything else first; BuzzSuite manages processing power automatically.

  \item \textbf{(Optional)} If you want to filter out unrealistic flight speeds, use the
        \textbf{Speed filter} box (Min/Max, or leave \textbf{``Auto Detect''} ticked --- it is by
        default) and click \textbf{``Run Speed Filter''}. This is usually not needed.
\end{enumerate}

That's it for this tab --- once tracking finishes, BuzzSuite automatically saves and combines the
data for you. There is no separate ``Concatenate'' or ``Load Data'' step anymore.

\subsection*{4.2 --- Or use the Experiment Dashboard}
\textit{(Tab: \texttt{2} \textperiodcentered\ \texttt{Analysis} $\to$ sub-tab \texttt{Multiple experiments})}

This does the same tracking job as 4.1, with a clearer per-experiment progress view. Use it for
one experiment, or several at once --- it works the same way either way:

\begin{enumerate}[leftmargin=*]
  \item Pick the experiment from the \textbf{Experiment} dropdown (it lists your 10 most recently
        opened experiments), or click \textbf{Browse\ldots} to pick any other one.

  \item Click \textbf{``Add to dashboard''}. This adds it as a row \textbf{without} starting
        tracking yet. If you want to track more than one experiment together, repeat this step
        for each of them first.

  \item Click \textbf{``Start''} on that row to begin tracking it --- or, if you've added more
        than one, click \textbf{``Start all queued''} to begin all of them together.

  \item Each row shows its own progress bar and status (\emph{queued} $\to$ \emph{starting}
        $\to$ \emph{tracking i/N segments} $\to$ \emph{concatenating} $\to$ \emph{done}), and a
        \textbf{Cancel} button if you need to stop it partway through.

  \item On the right, a \textbf{``Live file progress''} panel shows the progress of each
        individual video segment currently being tracked.
\end{enumerate}

If you do track more than one experiment at once, BuzzSuite automatically splits processing power
between whichever jobs are running --- there's no need to set worker numbers by hand or open a
second copy of the app. (You still need to finish Step~3, background \& borders, for each
experiment individually before tracking it, by switching which experiment is loaded on the
\texttt{1} \textperiodcentered\ \texttt{Setup} tab.)

% ----------------------------------------------------------------------
\section{Step 5 --- Make the activity plot}
\textit{(Tab: \texttt{3} \textperiodcentered\ \texttt{Plotting} $\to$ sub-tab \texttt{Activity (in-cage)})}

\begin{enumerate}[leftmargin=*]
  \item Click \textbf{``Add experiment''} (pick it from the recent-experiments dropdown, or
        \textbf{Browse\ldots}) to load its tracked data in.

  \item In the \textbf{dates table} below, tick which dates to include, and mark each one
        \textbf{LD} (normal light:dark cycle) or \textbf{DD} (constant dark) --- \textbf{LD} is
        the default for every date. Use \textbf{``All LD''} / \textbf{``All DD''} to set them all
        at once. Click \textbf{``Save LD/DD tags''} to keep this for next time.

  \item Under \textbf{Plot}, leave \textbf{Variable} on its default
        (\texttt{total\_pixels\_moved} --- overall movement per minute), or pick a different one
        from the dropdown.

  \item Choose how you want it laid out (\textbf{Output}):
  \begin{itemize}
    \item \textbf{Individual} --- one panel per date \emph{(default)}
    \item \textbf{Overlay} --- every date on one 24-hour axis, useful for comparing days
    \item \textbf{Continuous} --- one panel with real calendar time running along the bottom
  \end{itemize}

  \item Click \textbf{``Plot''} --- this is the \textbf{activity plot}. Click \textbf{``Save
        PNG''} to save it as an image file.
\end{enumerate}

\end{document}
